% ID: SEL-D0061-OB-E04
% ESTADO: cotejada. Segunda pasada visual y borrador compilado revisados.
% Fuente: Selectividad/sel_2010_matematicas_aplicadas/mat_apli/mat_apli_junio/mat_apli_junio/EXAMEN-ANDALUCÍA-5.pdf
% SHA-256: 5ea9fab18d8185843726e08994f39cf34f9b7dab32d86c310902a2f203dace2a
% Página: 2; opción B; ejercicio original 4.

Un agricultor piensa que la producción media por naranjo, en su finca, es de 88 kg o más. Para confirmar su creencia selecciona, al azar, 10 de sus naranjos, pesa su producción y obtiene como resultado, en kg, para cada uno de ellos:
\[
80,\quad83,\quad87,\quad95,\quad86,\quad92,\quad85,\quad83,\quad84,\quad95.
\]
Se acepta que la producción de un naranjo sigue una distribución Normal con desviación típica 5 kg.
\begin{enumerate}[label=\alph*)]
\item \textbf{(1.5 puntos)} Plantee el contraste de hipótesis unilateral que responda a las condiciones del problema y determine la región crítica para un nivel de significación $\alpha=0.05$.
\item \textbf{(1 punto)} Con los datos de esta muestra, \textquestiondown{}qué conclusión debe obtener el agricultor sobre la producción media por naranjo de su finca, utilizando ese mismo nivel de significación?
\end{enumerate}
